% =====================================================================
\chapter{IA generativa para 3D y validación de la geometría}
\label{tema:6}
% =====================================================================

Los modelos generativos producen ya objetos 3D a partir de un texto o de una imagen. Para un optimizador geométrico son a la vez una herramienta (generan candidatos o puntos de partida) y un riesgo (sus salidas pueden parecer correctas en una imagen y ser inservibles para fabricar o para un motor). Este tema describe tres familias de generadores con sus artículos originales: la destilación de un modelo de difusión 2D en DreamFusion~\cite{dreamfusion}, los generadores directos de funciones implícitas como Shap-E~\cite{shape} y los generadores autorregresivos de mallas como MeshGPT~\cite{meshgpt}. Después formaliza la validación de la geometría para fabricación y para motores, y la integra en un bucle de generación y verificación. Se apoya en la tarea~4 del curso 16-825 de Carnegie Mellon~\cite{cmu16825} y en el filtro de fabricación aditiva de Langelaar~\cite{langelaar}.

\section{Tres maneras de generar geometría}

\begin{definicion}[Muestreo por destilación de puntuación]\label{def:sds}
Sea $x=g(\theta)$ una imagen renderizada desde una vista aleatoria de una representación 3D con parámetros $\theta$ (en DreamFusion, un campo de radiancia), y sea $\hat\epsilon_\phi(z_t;y,t)$ la predicción del ruido de un modelo de difusión de imágenes preentrenado y congelado, condicionado por el texto $y$, sobre la versión ruidosa $z_t$ de $x$ en el instante $t$. El \emph{muestreo por destilación de puntuación} (SDS) actualiza $\theta$ con el gradiente~\cite[§2.1]{dreamfusion}
\[
\nabla_\theta\mathcal{L}_{\text{SDS}}=\E_{t,\epsilon}\Bigl[w(t)\bigl(\hat\epsilon_\phi(z_t;y,t)-\epsilon\bigr)\frac{\partial x}{\partial\theta}\Bigr],
\]
donde $\epsilon$ es el ruido añadido y $w(t)$ una ponderación.
\end{definicion}

La idea es que el modelo de difusión «sabe» cómo debería parecer una imagen del objeto descrito, y la diferencia entre el ruido que predice y el que realmente se añadió señala en qué dirección corregir la imagen; la regla de la cadena a través del renderizador derivable (\cref{tema:4}) lleva esa corrección a los parámetros 3D. DreamFusion usa el modelo base de $64\times 64$ de Imagen sin modificarlo y un campo de tipo mip-NeRF~360, y encuentra que hace falta una guía sin clasificador muy intensa~\cite{dreamfusion}. En la práctica se mezclan la predicción condicionada y la no condicionada, $\hat\epsilon=\hat\epsilon_u+s(\hat\epsilon_c-\hat\epsilon_u)$, con $s$ grande~\cite{cmu16825}. El método no necesita datos 3D, pero cada objeto es una optimización completa.

Los \emph{generadores directos} aprenden de pares de objetos 3D y textos. Shap-E entrena primero un codificador que transforma cada objeto en los parámetros de una función implícita, que puede renderizarse como malla texturizada o como campo de radiancia, y después un modelo de difusión condicionado sobre esos parámetros; genera objetos en segundos~\cite{shape}. Los \emph{generadores autorregresivos de mallas} tratan la malla como un texto: MeshGPT aprende un vocabulario de \emph{embeddings} geométricos de caras con un codificador de grafos y cuantización vectorial residual, y un transformer de solo decodificador predice el siguiente índice del vocabulario, como un LLM predice el siguiente token (módulos~12 y~15)~\cite{meshgpt}. El resultado son mallas compactas, parecidas a las que haría un artista, a diferencia de las mallas densas que se obtienen extrayendo la isosuperficie de un campo; a cambio, el tamaño está limitado por el contexto (los autores entrenan con mallas de menos de 800 caras).

Ninguna de las tres familias garantiza que el resultado sea una malla cerrada, orientada, sin autointersecciones, con la escala correcta y fabricable. Esas propiedades hay que comprobarlas.

\section{Validación de la geometría}

\begin{definicion}[Comprobaciones topológicas]\label{def:validacion}
Sea $(V,F)$ una malla triangular. Se cuentan las \emph{aristas dirigidas} $(a,b)$ de todas las caras (cada cara $(a,b,c)$ aporta $(a,b)$, $(b,c)$ y $(c,a)$) y las aristas no dirigidas $\{a,b\}$. La malla
\begin{enumerate}
  \item tiene \emph{borde} si alguna arista no dirigida pertenece a una sola cara;
  \item es \emph{no variedad} si alguna arista pertenece a tres o más caras;
  \item está \emph{mal orientada} si alguna arista dirigida aparece más de una vez.
\end{enumerate}
Una malla cerrada, variedad y orientada de forma coherente es \emph{estanca}.
\end{definicion}

\begin{proposicion}[Una cara invertida]\label{prop:cara-girada}
Sea una malla estanca e invértase la orientación de una única cara $f$ cuyas tres aristas son compartidas con caras distintas de $f$. Entonces la malla sigue cerrada y sin aristas no variedad, exactamente tres aristas dirigidas aparecen dos veces, y el volumen de la \cref{prop:volumen} disminuye en $\tfrac13\det(v_a,v_b,v_c)$, el doble de la contribución original de $f$.
\end{proposicion}

Por eso una cara invertida es un defecto silencioso: no abre agujeros y una comprobación de estanqueidad basada solo en recuentos de aristas no dirigidas no la detecta, pero estropea el volumen, la iluminación (la normal apunta hacia dentro) y la clasificación interior/exterior que usan los programas de laminado.

Para \emph{fabricación aditiva} se añaden comprobaciones geométricas:
\begin{itemize}[leftmargin=1.4em]
  \item \emph{Voladizos}: con dirección de fabricación $b$ unitaria, una cara de normal $n$ que mira hacia abajo necesita soporte si forma con la horizontal un ángulo menor que un umbral, habitualmente $45^\circ$; equivale a $n\cdot b<-\cos 45^\circ\approx-0{,}707$. Las caras apoyadas en la bandeja se excluyen. El filtro de Langelaar incorpora una restricción de este tipo dentro de la optimización de topología~\cite{langelaar}.
  \item \emph{Espesor mínimo}: con la SDF de la pieza (\cref{tema:1}), el espesor local es aproximadamente el doble de la distancia al borde en el eje medial; en la práctica se lanzan rayos desde la superficie hacia dentro y se mide la distancia a la cara opuesta.
  \item \emph{Volumen de impresión, unidades y escala}: la caja envolvente debe caber en la máquina y estar en milímetros.
  \item \emph{Cavidades cerradas}: en procesos de lecho de polvo o resina atrapan material sin consolidar.
\end{itemize}
Para un \emph{motor de videojuegos}: número de triángulos y de materiales dentro del presupuesto (\cref{tema:5}), coordenadas UV presentes y sin solapes, normales y tangentes coherentes, escala y punto de pivote convenidos, niveles de detalle generados y nombres según la convención del proyecto.

\section{Generar y verificar}

El esquema robusto es un bucle: el generador propone, el validador comprueba, un reparador corrige lo automático (soldar vértices duplicados, reorientar caras, cerrar agujeros pequeños) y el resultado se vuelve a validar. Si cada candidato pasa la validación con probabilidad $p$, generar varios y quedarse con el primero válido es eficaz (\cref{fig:bucle}).

\addtocounter{figure}{1}%

\begin{proposicion}[Lo mejor de $N$]\label{prop:mejorN}
Si $N$ candidatos superan la validación de forma independiente con probabilidad $p\in(0,1)$, la probabilidad de obtener al menos uno válido es $1-(1-p)^N$, y para alcanzar una probabilidad $1-\delta$ basta con
\[
N\ge\frac{\ln\delta}{\ln(1-p)} .
\]
\end{proposicion}

El validador se convierte así en una función de recompensa verificable, análoga a las recompensas verificables con que se ajustan los LLM (módulo~21, §6.2). Y hereda su riesgo: si el generador se ajusta para pasar el validador, aprenderá a satisfacer las comprobaciones sin satisfacer la intención (por ejemplo, mallas estancas que se autointersecan si no se comprueban las autointersecciones). Es la ley de Goodhart que el módulo~21 (§6.4) estudia como sobreoptimización de la recompensa. La defensa es la misma: comprobaciones independientes que no se usan para ajustar el generador y revisión humana por muestreo.

\section{Caso de estudio: validación de mallas generadas}

\addtocounter{definicion}{1}%

\begin{table}[h]
\centering\small
\begin{tabularx}{\textwidth}{@{}l>{\raggedright\arraybackslash}X>{\raggedright\arraybackslash}p{2.6cm}>{\raggedright\arraybackslash}X@{}}
\toprule
Herramienta & Para qué sirve & Licencia o tipo & Cuándo elegirla \\
\midrule
trimesh & estanqueidad, orientación, volumen, reparaciones simples & MIT & validador ligero en Python \\
PyMeshLab & filtros de reparación (no variedad, agujeros) & GPL-3.0 & reparación por lotes \\
Open3D & métricas geométricas y distancias entre mallas y nubes & MIT & comparar con una referencia \\
Blender (\texttt{bpy}) & comprobaciones y correcciones dentro del pipeline del artista & GPL-3.0 & validar recursos antes de exportarlos al motor \\
Shap-E & generación de funciones implícitas a partir de texto o imagen & MIT (pesos publicados) & prototipos y generación local reproducible \\
\bottomrule
\end{tabularx}
\caption{Tecnologías del tema~6.}
\label{tab:tec6}
\end{table}

\addtocounter{definicion}{1}%

